% - Repozitář Rust překladače
%    - Konkrétní use-case
%    - Chceme sledovat `rust-lang/rust`, který má spoustu PR, je integrovaný s team API, popsat jednoduše API, jaké se tam používají labely a k čemu.

\chapter{Repozitář Rust překladače}
\label{chap:rust_repository}

Rust je systémový programovací jazyk zaměřený na bezpečnost paměti a~výkon, jehož vývoj zastřešuje nezisková organizace Rust Foundation. Vývoj probíhá otevřeně na GitHubu v~rámci organizace rust-lang, která sdružuje stovky přispěvatelů a~spravuje desítky repozitářů. Samotný překladač je vyvíjen v~repozitáři rust-lang/rust, kam denně přispívají vývojáři z~celého světa s~různými rolemi a~oprávněními. Tato komplexita si vyžaduje automatizaci procesů, jako jsou přiřazování reviewerů, správa štítků nebo řízení merge fronty~\cite{rust_foundation,rust_lang_org}.

Tato kapitola se zaměřuje na specifika a~charakteristiku repozitáře Rust překladače. Z~důvodu jeho komplexity a~interního workflow je tento repozitář ideálním kandidátem pro analýzu a~sledování vývoje. Denně jsou zde otevírány a~uzavírány desítky \PR a~issues s~různými návrhy a~úpravami. Jsou zde rovněž popsána specifika relevantní pro co nejefektivnější vhled do stavu a~průběhu vývoje projektu.

\section{Rust Team \API}
\label{sec:rust_team_api}

Organizace rust-lang je rozdělena do specializovaných týmů, z~nichž každý odpovídá za konkrétní oblast projektu, například za vývoj překladače, návrh jazyka nebo správu infrastruktury. Členové týmů mají různá oprávnění, jako je schvalování \PR nebo přidělování štítků, přičemž tato oprávnění jsou svázána právě s~příslušností k~danému týmu.

Organizace rust-lang spravuje přístupová práva uživatelů prostřednictvím speciálního repozitáře \texttt{rust-lang/team}\footnote{\url{https://github.com/rust-lang/team}}, který obsahuje seznam týmů, jejich členů a~přidělených rolí. Tento repozitář funguje jako jediný zdroj pravdy (\textit{single source of truth}) pro informace o~tom, kdo má jaká oprávnění a~jaké týmy v~organizaci existují.

Struktura týmů a~jejich přiřazení uživatelů je definována pomocí souborů ve formátu TOML, které jsou spravovány prostřednictvím \PR. Na základě tohoto repozitáře je generováno veřejné Team \API využívané automatizačními boty v~repozitářích organizace.

Toto řešení umožňuje provoz automatizačních botů, jako je bors nebo triage, v~repozitáři překladače Rustu. Tito boti zajišťují přidělování kontroly \PR a~issues podle interních algoritmů a~automaticky přiřazují štítky označující aktuální stav daného \PR nebo issue.

\subsection{Přispěvatelé a~členové týmů}
\label{subsec:rust_team_api:contributors_members}

V~kontextu repozitáře rust-lang/rust je nutné rozlišovat dva pojmy.

\emph{Přispěvatel} (contributor) je jakýkoliv GitHub uživatel, který vytvořil \PR nebo issue v~repozitáři. Přispěvatel nemusí být přiřazen k~žádnému týmu. Přispěvatelé jsou identifikováni svým GitHub uživatelským jménem a~jednoznačným číselným identifikátorem.

\emph{Člen týmu} je přispěvatel, který je zároveň evidován v~Team \API jako člen nějakého týmu organizace rust-lang. Členové týmů mají zpravidla udělena dodatečná oprávnění odpovídající oblasti, za níž jejich tým odpovídá. Například právo schvalovat \PR nebo přiřazovat štítky. Jeden uživatel může být členem více týmů současně.

\subsection{Bezpečnost}
\label{subsec:rust_team_api:security}

Bezpečnost Team \API zajišťuje několik opatření. Prvním z~nich je zabránění neoprávněnému přístupu prostřednictvím spuštění škodlivého kódu v~\CI. Úspěšný útok by mohl vést k~získání plných práv k~celé infrastruktuře organizace. Ochrana spočívá v~tom, že nikdo kromě hlavních administrátorů, ani běžní správci (maintaineři), nemá přístup k~tajným tokenům, nemůže upravovat soubory pipeline, provádět sloučení neschváleného kódu ani si zvyšovat vlastní oprávnění.

Druhá skupina bezpečnostních opatření chrání před hrozbami manipulace s~obsahem repozitáře. Takový útok by mohl provést maintainer a~mohl by zahrnovat například změnu \URL domovské stránky za falešnou, eskalaci vlastních práv nebo práv napříč více týmy, případně smazání ostatních uživatelů. Ochranou proti těmto útokům je povinnost schválení změn ostatními maintainery a~zavedení kontrol v~\CI, které ověřují, zda změny v~souborech TOML nevedly k~podezřelé aktivitě~\cite{rust_team_docs}.

\section{Štítky}
\label{sec:rust_labels}

Výchozí stavy, které poskytuje GitHub (otevřeno, zavřeno, sloučeno), jsou pro mnoho malých projektů dostačující, avšak v~repozitáři překladače Rustu jsou nedostatečné. Proto jsou zde automaticky přiřazovány různé druhy štítků, jako například \code{S-waiting-for-author} nebo \code{S-waiting-for-review}, díky nimž je zřetelně patrné, v~jakém stavu se daný \PR nebo issue právě nachází.

Štítky lze primárně rozdělit do následujících kategorií:

\begin{itemize}
  \item \code{S-*} -- stavy, ve kterých se daný \PR nebo issue právě nachází. Příkladem jsou štítky \code{S-waiting-for-review} (\PR čeká na code review) nebo \code{S-waiting-for-author} (\PR čeká na úpravy od autora). Tyto štítky jsou zásadní pro sledování stavu v~průběhu vývoje.
  \item \code{T-*} -- štítky označující tým, jehož se daný \PR nebo issue týká. Například \code{T-lang} pro jazykový tým nebo \code{T-compiler} pro tým pracující na překladači.
  \item \code{A-*} -- štítky pro oblasti, kterých se daný \PR nebo issue dotýká. Například \code{A-macros} pro makra nebo \code{A-diagnostics} pro zprávy o~chybách a~upozorněních.
  \item \code{E-*} -- výzvy ke spolupráci (\textit{calls for participation}). Například \code{E-easy} označuje issue vhodnou pro nové přispěvatele a \code{E-needs-mcve} žádá o~minimalizaci reprodukce chyby~\cite{rust_issue_triaging}.
\end{itemize}

Typů štítků je v~repozitáři výrazně více, avšak výše uvedené kategorie jsou základní pro pochopení fungování systému štítků a~jeho možností pro sledování stavu a~kategorizaci jednotlivých \PR a~issues~\cite{rust_issue_triaging}.

\section{Proces sloučení pull requestů}
\label{sec:rust_merge_process}

Sloučení \PR v~repozitáři rust-lang/rust neprobíhá standardním způsobem pomocí merge na GitHubu. Místo toho je tento proces řízen specializovaným botem bors, který zajišťuje, že každý \PR je před sloučením otestován v~aktuálním stavu větve \code{main}. Tím je zabráněno situaci, kdy \PR sice prochází testy v~okamžiku code review, ale po sloučení s~mezitím aktualizovanou větví \code{main} testy selžou~\cite{rust_lang_bors}.

\subsection{Bot bors}
\label{subsec:rust_merge_process:bors}

Bors je bot pro správu fronty sloučení (\textit{merge queue}). Jeho instance pro repozitář rust-lang/rust je dostupná na adrese \url{https://bors.rust-lang.org/queue/rust} a~vystupuje pod účtem \texttt{@bors} na GitHubu. Bot je naprogramován v~jazyce Rust a~jeho zdrojový kód je dostupný v~repozitáři \texttt{rust-lang/bors}~\cite{rust_lang_bors,forge_rust_bors}.

\subsection{Průběh sloučení pull requestu}
\label{subsec:rust_merge_process:workflow}

Průběh sloučení \PR v~repozitáři rust-lang/rust lze shrnout do následujících kroků~\cite{rustc_contributing}:

\begin{enumerate}
  \item \textbf{Otevření \PR} -- Přispěvatel otevře \PR s~navrhovanými změnami. Bot \texttt{@rustbot} automaticky přiřadí reviewera na základě změněných souborů a~přidělí odpovídající štítky.
  \item \textbf{Code review} -- Přiřazený reviewer přezkoumá změny a~případně si vyžádá úpravy (štítek \code{S-waiting-on-author}). Jakmile je \PR připraven ke sloučení, reviewer jej schválí příkazem \code{@bors r+} v~komentáři.
  \item \textbf{Fronta sloučení} -- Příkaz \code{r+} zařadí \PR do fronty sloučení bota bors.
  \item \textbf{Spuštění \CI} -- Bors vytvoří pomocnou větev \code{automation/bors/auto}, do které sloučí změny z~\PR s~aktuálním stavem větve \code{main}, a~spustí kompletní sadu testů \CI.
  \item \textbf{Sloučení} -- Pokud testy projdou, bors sloučí \PR do větve \code{main} a~uzavře jej. Pokud testy selžou, bors \PR zamítne a~informuje autora o~výsledku.
\end{enumerate}
